% 第一章：绪论
% 论文第一章一般交代研究背景、现状与本文工作。
% 本章演示了常用的章节层级与图表引用方式。
\chapter{绪\quad 论}

\section{研究背景与意义}
% 用 1～2 段说明课题来源、要解决的实际问题与研究意义
在此阐述研究背景与课题意义。可以先介绍问题的实际背景，
再指出现有方法存在的不足，从而引出本文的研究动机。

\section{国内外研究现状}
\subsection{方向一的研究现状}
按技术路线或时间顺序综述代表性工作，引用示例：统计学习理论
框架见文献~\cite{vapnik2000}，深度学习方法的综述见
文献~\cite{lecun2015}，Transformer 架构由文献~\cite{vaswani2017} 提出。

\subsection{方向二的研究现状}
继续综述与本课题相关的其他方向，并总结现有工作的不足。

\section{主要研究内容}
% 用一段话 + 条目列出本文的三到四项工作
本文的主要研究内容包括：
\begin{itemize}
  \item 研究内容一：对问题进行建模与形式化描述；
  \item 研究内容二：设计并实现示例系统；
  \item 研究内容三：通过实验验证方法的有效性。
\end{itemize}

\section{论文组织结构}
% 逐章一句话说明全文结构
本文共分为三章：第一章介绍研究背景与主要内容；
第二章介绍方法设计与实验分析；第三章总结全文并展望未来工作。
